\section{Independence of finite divisorial coordinate presentations}

\begin{lemma}[Reindexing the original tuples does not change their Jacobian ideal]\label{lem:sc0-dvr-jacobian-finite-reindex}
\lean{EnvelopeProperty.Char0.dvrOriginalSourceJacobianIdeal_reindex}
\leanok
\uses{lem:sc0-original-dvr-tuple-field-jacobian}
Fix an actual uniformizer and lifted finite residue-transcendence basis of an
actual divisorial DVR $V$, and a source subalgebra $C\subseteq V$. Any two
finite enumerations of those differential coordinates give the same ideal
generated by determinants of all original source-function tuples.
\end{lemma}
\begin{proof}
For enumerations $e,e'$, replace an original tuple $w$ by
$w'=w\circ e^{-1}\circ e'$. After embedding into the actual fraction field,
the two determinant formulas are identical: both are the Jacobian of the
tuple indexed by the differential coordinates. The DVR embedding is
injective, so the original DVR determinants are equal. This replacement is
invertible, using the reverse composite. Each set of generators is therefore
contained in the other generated ideal, which proves equality.
\end{proof}

\begin{theorem}[All finite divisorial coordinate presentations give the same actual ideal]\label{thm:sc0-dvr-jacobian-presentation-independent}
\lean{EnvelopeProperty.Char0.dvrOriginalSourceJacobianIdeal_presentation_independent}
\leanok
\uses{lem:sc0-dvr-jacobian-finite-reindex,thm:sc0-original-dvr-jacobian-ideal-coordinate-independent}
For a fixed actual essentially finite-type divisorial DVR and original
source subalgebra, its genuine all-original-source Jacobian ideal is
independent of the finite number of lifted residue coordinates, their
values, the uniformizer, and the finite enumeration of differential
coordinates.
\end{theorem}
\begin{proof}
Each finite coordinate presentation gives a basis of the differentials of
the same actual fraction field. The dimension of that vector space equals
the cardinality of either basis. Hence the two total numbers of coordinates
are equal; subtracting the one uniformizer coordinate shows that the two
numbers of residue coordinates are equal. Transport along those equalities.
The fixed-index coordinate-invariance theorem identifies the ideals for the
two actual uniformizer and residue-coordinate choices. The preceding
reindexing lemma then removes the difference in enumerations.
\end{proof}
